\section{Conclusion}
\label{sec:conclusion}

We pre-registered an audit of one prominent RAG faithfulness claim
--- the refinement paradox --- and used it to expose how
evaluation choices shape the conclusions that get reported. The
published $n{=}30$ paradox magnitude collapsed at $n{=}500 \times 5$
seeds; the proposed causal mechanism (CCS) was not supported by a
matched-similarity intervention; the choice of faithfulness scorer
moved the magnitude by an order of magnitude on the same paired
outputs; the published threshold $\tau$ did not transfer cleanly
across datasets; head-to-head comparison against published
baselines showed no method dominates.

A structural finding survived: at matched per-passage similarity,
coherence-preserving uninformative noise produces a smaller
faithfulness drop than random off-topic noise. Coherence is real
in some operationalization, even if the published one (CCS) is
not the right one.

We release the benchmark, scripts, pre-registrations, and
per-query CSVs at
\href{https://github.com/Saket-Maganti/rag-hallucination-detection}{github.com/Saket-Maganti/rag-hallucination-detection}
and DOI \href{https://doi.org/10.5281/zenodo.19757291}{10.5281/zenodo.19757291},
on free hardware, so the same audit can be run on any candidate
RAG faithfulness intervention before its claims are finalized.
We argue that pre-registration, multi-scorer reporting, and
explicit cross-dataset $\tau$ disclosure should become the minimum
bar for empirical claims in RAG faithfulness research.

\paragraph{Acknowledgments.}
\ifanonymous
We omit acknowledgments under double-blind review.
\else
We thank the prior reviewers whose feedback prompted this audit,
and the maintainers of Hugging Face Hub, Zenodo, Kaggle, and
Ollama for the free infrastructure that made the work possible
under a zero-dollar constraint.
\fi
